% This is samplepaper.tex, a sample chapter demonstrating the
% LLNCS macro package for Springer Computer Science proceedings;
% Version 2.21 of 2022/01/12
%
\documentclass[runningheads]{llncs}
%
\usepackage[T1]{fontenc}
% T1 fonts will be used to generate the final print and online PDFs,
% so please use T1 fonts in your manuscript whenever possible.
% Other font encondings may result in incorrect characters.
%
\usepackage{graphicx}
% Used for displaying a sample figure. If possible, figure files should
% be included in EPS format.
%
% Formatting packages
\usepackage{csquotes}
%
% Layout packages
\usepackage[margin=1in]{geometry}
%
% If you use the hyperref package, please uncomment the following two lines
% to display URLs in blue roman font according to Springer's eBook style:
%\usepackage{color}
%\renewcommand\UrlFont{\color{blue}\rmfamily}
%\urlstyle{rm}
%
\begin{document}
%
\title{Sustainable Resource Allocation in Urban School Settings
(STRAUSS)}
%
\titlerunning{STRAUSS}
% If the paper title is too long for the running head, you can set
% an abbreviated paper title here
%
\author{Andrew Poile\inst{1}\orcidID{0000-1111-2222-3333} \and
Baharak Rastegari\inst{1}\orcidID{0000-0002-0985-573X} \and
Vahid Yazdanpanah\inst{1}\orcidID{0000-0002-4468-6193} \and
Sebastian Stein\inst{1}\orcidID{0000-0003-2858-88576}}
%
\authorrunning{A. Poile et al.}
% First names are abbreviated in the running head.
% If there are more than two authors, 'et al.' is used.
%
\institute{
    University of Southampton, Southampton SO17 1BJ, UK
%   \email{lncs@springer.com}\\
%   \url{http://www.springer.com/gp/computer-science/lncs}
    }
%
\maketitle              % typeset the header of the contribution
%
\begin{abstract}
In the 1988 Education Reform Act \cite{participationEducationReformAct}, the UK introduced school choice mechanisms into England and Wales. 

\keywords{Matching \and School Choice \and Mechanism Design.}
\end{abstract}
%
%
%
\section{Introduction}
\subsection{School Choice (change later)}
Empirical evidence suggests that school choice has increased school segregation in many cities and contexts \cite{wilsonSchoolChoiceCity2019}. Particularly in England, the geographical constraints of the system with respect to over-subscription criteria, which either implicitly or explicitly increase the priority of students proximal to the school, restrict the feasible choices of disadvantaged families \cite{burgessParentalChoicePrimary2011a,burgessSchoolChoiceEngland2019}, and result in socio-economic segregation in schools. This restriction in choice has the effect of re-introducing the geographical inequities that the school choice system was supposed to eliminate.

\section{Related Literature}

\section{Model}
\subsection{Components}
\begin{itemize}
    \item 3 disjoint sets; students, colleges, routes:
    \subitem $S = \{s_1,...,s_{n_S}\}, C = \{c_1,...,c_{n_C}\}, R = \{r_1,...,r_{n_R}\}$.
    \item 2 sets constructed from elements of the 3 sets above; route-college tuples and route-student tuples:
    \subitem $A = \{a_{11},...,a_{n_R n_C}\}$ where eg. for student $s_1: a_{11} = (r_1,c_1) \succ ... \succ a_{23} = (r_2,c_3) \succ ...$ etc,
    \subitem $B = \{b_{11},...,b_{n_R n_S}\}$ where eg. for college $c_1: b_{11} = (r_1,s_1) \succ ...$, for $c_3: b_{21} = (r_2,s_1) \succ ...$, etc.
    \item College over-subscription priorities divided into 2 priority brackets given by:
    \subitem Students \textit{within} a certain distance of a school selecting that school \textbf{without} transport AND students \textit{within} a selected area served by a route, selecting a school \textbf{with} transport are placed in the same priority bracket; all other students placed in the bracket below this.
    \item Routes can have a capacity of more than 1: $q_r \geq 1$.
    \item Each route serves one college. Routes are effectively taxi services.
    \item Routes are only available in selected areas (only feasible for a subset of students $S_{r_i} \in S$ for a route i). The area would be selected based on SES indicators, such as eligibility for free school meals, to maximise the impact of improving priority for students (current model selects a random subset of students to be served by routes, and each student is served by one route).
    \item Eligible students (those served by routes) can select either a college, or route-college tuple, eg. in student $s_1$'s preference list: $a_{11} = (r_1,c_1) \succ a_{23} = (r_2,c_3) \succ c_1 \succ c_3 \succ ...$ etc. Each object $a$ contains exactly one college. Ineligible students may only select colleges and are unable to select route-college tuples.
    \item Each route-college tuple has a \textbf{corresponding route-student tuple}, where the route is shared between the two tuples, eg. for $c_1: b_{11} = (r_1,s_1), s_1, ...$, and for $c_3: b_{21} = (r_2,s_1), s_1, ...$, etc from the route-college tuples above. All route-college tuples have their corresponding route-student tuples added to the top priority bracket in the priority list of the college in the route-college tuple.
    \item Student preferences over colleges and route-college tuples (that are available to them), and college priorities over students and route-student tuples are strict and complete.
\end{itemize}

\section{Results}

% \subsubsection{Sample Heading (Third Level)} Only two levels of
% headings should be numbered. Lower level headings remain unnumbered;
% they are formatted as run-in headings.

% \paragraph{Sample Heading (Fourth Level)}
% The contribution should contain no more than four levels of
% headings. Table~\ref{tab1} gives a summary of all heading levels.

% \begin{table}
% \caption{Table captions should be placed above the
% tables.}\label{tab1}
% \begin{tabular}{|l|l|l|}
% \hline
% Heading level &  Example & Font size and style\\
% \hline
% Title (centered) &  {\Large\bfseries Lecture Notes} & 14 point, bold\\
% 1st-level heading &  {\large\bfseries 1 Introduction} & 12 point, bold\\
% 2nd-level heading & {\bfseries 2.1 Printing Area} & 10 point, bold\\
% 3rd-level heading & {\bfseries Run-in Heading in Bold.} Text follows & 10 point, bold\\
% 4th-level heading & {\itshape Lowest Level Heading.} Text follows & 10 point, italic\\
% \hline
% \end{tabular}
% \end{table}


% \noindent Displayed equations are centered and set on a separate
% line.
% \begin{equation}
% x + y = z
% \end{equation}
% Please try to avoid rasterized images for line-art diagrams and
% schemas. Whenever possible, use vector graphics instead (see
% Fig.~\ref{fig1}).

% \begin{figure}
% %\includegraphics[width=\textwidth]{fig1.eps}
% \caption{A figure caption is always placed below the illustration.
% Please note that short captions are centered, while long ones are
% justified by the macro package automatically.} \label{fig1}
% \end{figure}

\begin{theorem}
The modified mechanism terminates.
\end{theorem}

\begin{proof}
    \begin{enumerate}
        \item The number of agents, the length of the preference and priority lists, and the college and route capacities are all finite.
        \item Students propose to objects in order of preference, and if they are accepted by an object they will not propose again unless they are unmatched with that object at a later stage. If they are unmatched later, they will not propose to objects that they have been unmatched with, so their preference lists will be weakly decreasing with every iteration of the algorithm.
        \item Even if all students were unacceptable to all colleges, and every student was rejected at every stage, there would be a maximum of $n \times m$ proposals, where $n$ is the number of students and $m$ is the number of colleges.
    \end{enumerate}
\end{proof}

\begin{theorem}
The modified mechanism still returns a stable matching.
\end{theorem}
%
% the environments 'definition', 'lemma', 'proposition', 'corollary',
% 'remark', and 'example' are defined in the LLNCS documentclass as well.
%
\begin{proof}
    Suppose there exists some student $s$ and [college $c$ OR tuple $(r,c)$] which are not matched together in $\mu$ after termination of the algorithm such that neither $(s,c)$ nor $(s,(r,c))$ are in the matching $\mu$.

    Suppose also that $s$ prefers [$c$ OR $(r,c)$] to $\mu(s)$ and hence [$c$ OR $(r,c)$] must be acceptable to $s$, and $s$ must have applied to [$c$ OR $(r,c)$] before applying to $\mu(s)$. In the case where $\mu(s)$ is empty, the student $s$ must be unmatched, and hence must have exhausted their entire preference list before ending up unmatched.

    Since $s$ is not matched to [$c$ OR $(r,c)$] when the algorithm terminates:
    \begin{itemize}
        \item IF $s$ is NOT eligible for $r$: either $s$ applied to $c$, was accepted and then later rejected by $c$ in favour of a higher priority student/routed student, OR $c$ was already at capacity at the point of application and was rejected as the lowest priority student/routed student.
        \item ELSE IF $s$ IS eligible for $r$ and ONLY applied to $(r,c)$; either:
        \begin{itemize}
            \item $s$ applied to $(r,c)$, was accepted and was later rejected from $r$ (and since $r$ only serves $c$ is therefore also rejected from $c$) in favour of a higher priority student routed to $c$, OR $r$ was already at capacity at the point of application and $s$ was rejected as the lowest priority student routed to $r$,
            \item OR
            \item $s$ applied to $(r,c)$, was accepted and later rejected from $c$ in favour of a higher priority student/routed student, OR $c$ was already at capacity at the point of application and was rejected as the lowest priority student/routed student.
        \end{itemize}
        \item ELSE IF $s$ IS eligible for $r$ and applied to BOTH $(r,c)$ AND $c$; BOTH:
        \begin{itemize}
            \item $s$ applied to $(r,c)$, was accepted and was later rejected from $r$ (and therefore also $c$) in favour of a higher priority routed student, OR $r$ was already at capacity at the point of application and $s$ was rejected as the lowest priority routed student
            \item AND
            \item $s$ applied to $c$, was accepted and was later rejected by $c$ in favour of a higher priority student/routed student, OR $c$ was already at capacity at the point of application and was rejected as the lowest priority student/routed student.
        \end{itemize}
    \end{itemize}

    A key aspect of the stability of our algorithm is the order in which the oversubscription is resolved in the case where both college and route are oversubscribed: as long as route oversubscription is resolved first, college oversubscription is resolved simultaneously and the number of remaining seats of a college \textit{can only weakly decrease with each step of the algorithm.}

    Because of this, \textit{the lowest priority student in a college's preference list can \textbf{only} weakly improve.} Hence, the cases above hold not only at the point of rejection, but at the termination of the algorithm.

    Additionally, matchings may be stable, even if routes are undersubscribed.

    Below are examples that showcase the above statements:

    \subsection{Examples}
    \subsubsection{Case 1}
    Consider the case:

    \begin{center}
        \begin{tabular}{ c|l }
            College&Priorities \\
            \hline
            $c_1$&$(r_1,s_4),(r_1,s_5),(r_1,s_6),s_3,(r_1,s_7)$ \\
        \end{tabular}
    \end{center}

    where $c_1$ has a capacity of 3 and $r_1$ has a capacity of 2.

    If the students propose in the reverse order of college $c_1$'s priority list, then we see that:

    \begin{enumerate}
        \item $(r_1,s_6),s_3,(r_1,s_7)$ are initially assigned to $c_1$,
        \item $(r_1,s_5)$ applies and displaces $(r_1,s_7)$ since $r_1$ has reached capacity (as well as $c_1$) and $(r_1,s_7)$ is the lowest priority:
        \begin{itemize}
            \item $r_1$ and $c_1$ reach capacity simultaneously, but since the lowest priority item, $(r_1,s_7)$, is simultaneously the lowest priority item in $c_1$'s priority list \textit{and} the lowest priority routed student, the order in which the overcapacity problem is solved does not matter.
        \end{itemize}
        \item $(r_1,s_4)$ applies and now there is a problem: both $r_1$ and $c_1$ are already at capacity, and the order in which the oversubscription is resolved is important:
        \begin{itemize}
            \item If $c_1$'s oversubscription is resolved first, then $s_3$ will be rejected from the college, after which $(r_1,s_6)$ will be rejected from both route and college, resulting in the final allocation $c_1:(r_1,s_4),(r_1,s_5)$, leaving the college \textit{undersubscribed}, but importantly violating one of the principles that guarantees termination and stability: that colleges only \enquote{trade up.}
            \item On the contrary, if the route oversubscription is resolved first, then only $(r_1,s_6)$ will be rejected from both route and college, which satisfies both capacity conditions simultaneously, and the final allocation will be: $c_1:(r_1,s_4),(r_1,s_5),s_3$.
        \end{itemize}
    \end{enumerate}

    \subsubsection{Case 2}
    Consider the next case where the priority list is prepended with $s_1,s_2$, who are local students in the upper priority bracket.

    \begin{center}
        \begin{tabular}{ c|l }
            College&Priorities \\
            \hline
            $c_1$&$s_1,s_2,(r_1,s_4),(r_1,s_5),(r_1,s_6),s_3,(r_1,s_7)$ \\
    \end{tabular}
    \end{center}

    After following the same procedure as above, $s_1,s_2$ now apply to $c_1$ in reverse order as before, displacing first $s_3$, then $(r_1,s_5)$. This will leave $r_1$ \textit{undersubscribed}. This is not an issue here however because there are no students, routed or otherwise who have a higher priority than the assigned students, and that would be able to change the assignment, as all other acceptable students have already been rejected, and they cannot re-apply later.
\end{proof}

\begin{theorem}
    The modified mechanism is student optimal if run student-oriented. No student is rejected by an achievable object and so the algorithm returns the student optimal stable matching.
\end{theorem}

\begin{proof}
    Suppose the algorithm is at iteration $i$ and no student has been unmatched with an object that is achievable to them:
    \begin{enumerate}
        \item At this iteration $i$ suppose route $r$, which serves only college $c$ and inherits its priority list, is at capacity and student $s$ who applied to $c$ with route $r$ (henceforth called $(r,s)$), is rejected from both $r$ and $c$:
        \begin{itemize}
            \item If $(r,s)$ is rejected in favour of $(r,s')$, then $(r,s)$ is lower in $c$'s priority list than $(r,s')$.
            \item Since there have been no rejections up until this iteration, and students apply to their top preferences first in the algorithm, we know that $(r,s')$ prefers $c$ over all other colleges.
            \item If we consider a matching $\mu$ where $(r,s)$ is matched to $c$ and all other students are matched to achievable objects, we know it will be unstable since $(r,s')$ prefers $c$ to its current partner, and $c$ prefers $(r,s')$ to $(r,s)$.
            \item Hence, there is no stable matching where $(r,s)$ is matched to $c$, and so they are unachievable for each other.
        \end{itemize}
        \item Else at this iteration suppose college $c$ is at capacity and rejects student $s$ or $(r,s)$.
        \begin{itemize}
            \item If $s$ or $(r,s)$ is rejected in favour of $s'$ or $(r,s')$, then $s$ or $(r,s)$ is lower in $c$'s priority list than either of $s'$ or $(r,s')$.
            \item Clearly $s'$ or $(r,s')$ prefer $c$ over all other colleges since students apply to their top preferences first in the algorithm, and no student has been rejected as of iteration $i$.
            \item If we consider a matching $\mu'$ where $s$ or $(r,s)$ is matched to $c$ and all other students are matched to achievable objects, we know it will be unstable since $s'$ or $(r,s')$ prefers $c$ to its current partner, and $c$ prefers $s'$ or $(r,s')$ to $s$ or $(r,s)$.
            \item Hence, there is no stable matching where $s$ or $(r,s)$ is matched to $c$, and so they are unachievable for each other.
        \end{itemize}
    \end{enumerate}
\end{proof}
% For citations of references, we prefer the use of square brackets
% and consecutive numbers. Citations using labels or the author/year
% convention are also acceptable. The following bibliography provides
% a sample reference list with entries for journal
% articles~\cite{ref_article1}, an LNCS chapter~\cite{ref_lncs1}, a
% book~\cite{ref_book1}, proceedings without editors~\cite{ref_proc1},
% and a homepage~\cite{ref_url1}. Multiple citations are grouped
% \cite{ref_article1,ref_lncs1,ref_book1},
% \cite{ref_article1,ref_book1,ref_proc1,ref_url1}.

\begin{credits}
\subsubsection{\ackname} A bold run-in heading in small font size at the end of the paper is
used for general acknowledgments, for example: This study was funded
by X (grant number Y).

\subsubsection{\discintname}
It is now necessary to declare any competing interests or to specifically
state that the authors have no competing interests. Please place the
statement with a bold run-in heading in small font size beneath the
(optional) acknowledgments\footnote{If EquinOCS, our proceedings submission
system, is used, then the disclaimer can be provided directly in the system.},
for example: The authors have no competing interests to declare that are
relevant to the content of this article. Or: Author A has received research
grants from Company W. Author B has received a speaker honorarium from
Company X and owns stock in Company Y. Author C is a member of committee Z.
\end{credits}
%
% ---- Bibliography ----
%
% BibTeX users should specify bibliography style 'splncs04'.
% References will then be sorted and formatted in the correct style.
%
\bibliographystyle{splncs04}
\bibliography{MATCH-UP.bib}
%
% \begin{thebibliography}{8}
% \bibitem{ref_article1}
% Author, F.: Article title. Journal \textbf{2}(5), 99--110 (2016)

% \bibitem{ref_lncs1}
% Author, F., Author, S.: Title of a proceedings paper. In: Editor,
% F., Editor, S. (eds.) CONFERENCE 2016, LNCS, vol. 9999, pp. 1--13.
% Springer, Heidelberg (2016). \doi{10.10007/1234567890}

% \bibitem{ref_book1}
% Author, F., Author, S., Author, T.: Book title. 2nd edn. Publisher,
% Location (1999)

% \bibitem{ref_proc1}
% Author, A.-B.: Contribution title. In: 9th International Proceedings
% on Proceedings, pp. 1--2. Publisher, Location (2010)

% \bibitem{ref_url1}
% LNCS Homepage, \url{http://www.springer.com/lncs}, last accessed 2023/10/25
% \end{thebibliography}
\end{document}